\documentclass[11pt]{article}
\pdfinfoomitdate=1
\pdftrailerid{}
\pdfsuppressptexinfo=15
\usepackage[a4paper,margin=1in]{geometry}
\usepackage{amsmath,amssymb,amsthm,booktabs,microtype}
\usepackage[numbers,sort&compress]{natbib}
\usepackage{cleveref}

\newtheorem{theorem}{Theorem}[section]
\newtheorem{corollary}[theorem]{Corollary}
\newtheorem{proposition}[theorem]{Proposition}
\theoremstyle{remark}
\newtheorem{remark}[theorem]{Remark}
\newcommand{\cP}{\mathcal P}
\newcommand{\cA}{\mathcal A}

\title{A Double-Alias Parity-Phase Obstruction at the\
Critical Order}
\author{Bin Wang}
\date{August 2026}

\begin{document}
\maketitle

\begin{abstract}
RH-344 completes the physical boundary-orbit extraction at the first alias
and gives
\[
 p_{\sigma,k,2k}=Y_k+\cP_{\sigma,2k}-S_k,
 \qquad
 Y_k=\mathcal T_k^{\rm rest}-d_{\sigma,k,2k},
 \qquad
 S_k=\cA_{k,2k}+F_k^{\rm orb}.
\]
Here $S_k/\cA_{k,2k}\to2$ and $\cA_{k,2k}/H_k\to\infty$.  Combining this
with the actual RH-326 phase law
\[
 \frac{\cP_{\sigma,2k}}{\cA_{k,2k}}
 =C_*C_M\lambda^{\eta_\sigma}\{1+o(1)\}
\]
shifts the scalar balance from one alias packet to two.  The unique symbolic
phase is
\[
 \eta_2=\frac{\log(2/(C_*C_M))}{\log\lambda}.
\]
If the actual orbit-free remainder satisfies $Y_k=o(H_k)$, then every fixed
phase other than $\eta_2$ has a divergent critical weighted contribution.
This is a conditional physical obstruction, not aggregate nonclosure.

At $\eta_2$, the proved relative $o(1)$ phase law remains exponentially too
weak for the necessary relative precision $o((\beta R)^{-2k})$.  We prove an
exact scalar information-class theorem.  Two parity eigenvalue sequences of
the form $\lambda_{-,k}=-(1-\delta_k)$ both satisfy
$\delta_k=C_*\sqrt{\sigma_k}\{1+o(1)\}$ and the exact even parity-packet
formula.  One makes the critical scalar residual zero; the other leaves
$\cA_{k,2k}/k$, whose weighted contribution diverges.  These scalar sequences
are not two noisy operators.  Actual signed compensation, the strict prefix,
determinant gluing, and Gates A--E remain open.  No Riemann-hypothesis
conclusion follows.
\end{abstract}

\section{The complete critical scalar ledger}

Use the physical natural clock and target
\begin{equation}\label{eq:clock}
 k=\frac{\log(1/\sigma)}{2\log\lambda}+O(1),\qquad
 \eta_\sigma=k-\frac{\log(1/\sigma)}{2\log\lambda},\qquad
 H_k=kR^{-2k},\quad R=\frac75.
\end{equation}
RH-344 gives the exact direct critical coefficient
\begin{equation}\label{eq:critical}
 p_{\sigma,k,2k}=Y_k+\cP_{\sigma,2k}-S_k,
 \qquad
 Y_k=\mathcal T_k^{\rm rest}-d_{\sigma,k,2k},
 \qquad
 S_k=\cA_{k,2k}+F_k^{\rm orb}.
\end{equation}
The two positive deterministic packets obey
\begin{align}
 \cA_{k,2k}
 &=\frac{2k}{C_M}\beta^{2k}\{1+o(1)\},
 \label{eq:A-scale}\\
 F_k^{\rm orb}
 &=\frac{2k}{C_M}\beta^{2k}\{1+o(1)\},
 \label{eq:F-scale}
\end{align}
where $\beta=(r_H\sqrt\lambda)^{-1}$ and $\beta R>1$.  Therefore
\begin{align}
 \frac{S_k}{\cA_{k,2k}}&\longrightarrow2,
 \label{eq:S-over-A}\\
 S_k&=\frac{4k}{C_M}\beta^{2k}\{1+o(1)\},
 \qquad
 \frac{S_k}{H_k}=\frac4{C_M}(\beta R)^{2k}\{1+o(1)\}
 \longrightarrow\infty.
 \label{eq:S-scale}
\end{align}
These are exact physical deterministic constituents.  No assumption has yet
been made about the signed orbit-free remainder $Y_k$
\citep{WangCompleteOrbit2026}.

The actual parity eigenvalue satisfies the square-root boundary-layer law
\begin{equation}\label{eq:delta}
 \lambda_-(\sigma)=-(1-\delta_\sigma),\qquad
 \delta_\sigma=C_*\sqrt\sigma+o(\sqrt\sigma),\qquad C_*>0,
\end{equation}
and its even packet is
\begin{equation}\label{eq:P}
 \cP_{\sigma,2k}=r_H^{-2k}\{1-(1-\delta_\sigma)^{2k}\}.
\end{equation}
RH-326 proves, for bounded physical phase,
\begin{equation}\label{eq:phase-law}
 \frac{\cP_{\sigma,2k}}{\cA_{k,2k}}
 =C_*C_M\lambda^{\eta_\sigma}\{1+o(1)\}
\end{equation}
\citep{WangParityBoundary2026,WangFirstAlias2026}.

\section{Off-balance scalar obstruction}

Put
\begin{equation}\label{eq:gamma}
 \gamma(\eta)=C_*C_M\lambda^\eta.
\end{equation}
Since $\lambda>1$, $\gamma$ is strictly increasing.  The complete scalar
demand has the unique leading balance phase
\begin{equation}\label{eq:eta-two}
 \boxed{\eta_2=\frac{\log(2/(C_*C_M))}{\log\lambda}},
 \qquad \gamma(\eta_2)=2.
\end{equation}
This differs from the single-alias equation $\gamma=1$ because the complete
raw orbit contributes the second alias-sized packet.

\begin{theorem}[Conditional physical off-balance obstruction]
\label{thm:off-balance}
Consider an actual physical sequence with bounded phase
$\eta_\sigma\to\eta$.  If
\begin{equation}\label{eq:Y-small}
 Y_k=o(H_k)
\end{equation}
and $\eta\ne\eta_2$, then
\begin{equation}\label{eq:divergence}
 \frac{|p_{\sigma,k,2k}|}{2H_k}\longrightarrow\infty.
\end{equation}
More precisely,
\begin{equation}\label{eq:divergence-rate}
 \frac{|p_{\sigma,k,2k}|}{2H_k}
 =\frac{|\gamma(\eta)-2|}{C_M}
 (\beta R)^{2k}\{1+o(1)\}.
\end{equation}
\end{theorem}

\begin{proof}
By \eqref{eq:A-scale}, $H_k/\cA_{k,2k}\to0$, so
\eqref{eq:Y-small} gives $Y_k/\cA_{k,2k}\to0$.  Divide
\eqref{eq:critical} by $\cA_{k,2k}$ and use
\eqref{eq:S-over-A} and \eqref{eq:phase-law}:
\[
 \frac{p_{\sigma,k,2k}}{\cA_{k,2k}}
 \longrightarrow\gamma(\eta)-2\ne0.
\]
Finally \eqref{eq:A-scale} divided by $2H_k$ gives
$\cA_{k,2k}/(2H_k)=C_M^{-1}(\beta R)^{2k}\{1+o(1)\}$.
This proves both claims.
\end{proof}

\begin{remark}[Why this is not aggregate nonclosure]
The hypothesis \eqref{eq:Y-small} is a statement about the actual signed
orbit-free raw remainder and head defect.  No repository theorem proves it.
The theorem stops parity as a scalar-only compensation mechanism off balance;
it does not stop cancellation by $Y_k$.
\end{remark}

\section{The balance phase still misses target precision}

At $\eta=\eta_2$, the leading ratio in \eqref{eq:phase-law} equals two, but
critical closure requires much more.

\begin{proposition}[Necessary balance-phase precision]
\label{prop:precision}
Under $Y_k=o(H_k)$, critical closure is equivalent to
\begin{equation}\label{eq:scalar-match}
 \cP_{\sigma,2k}=S_k+o(H_k).
\end{equation}
Relative to $S_k$, the required precision is
\begin{equation}\label{eq:relative-precision}
 o\!\left(\frac{H_k}{S_k}\right)
 =o\!\left((\beta R)^{-2k}\right).
\end{equation}
The source law \eqref{eq:phase-law} supplies only relative $o(1)$ and hence
does not decide \eqref{eq:scalar-match}, even at $\eta_2$.
\end{proposition}

\begin{proof}
Rearrange \eqref{eq:critical} and use $Y_k=o(H_k)$.  Equation
\eqref{eq:S-scale} gives \eqref{eq:relative-precision}.  The final statement
is the strict distinction between an unspecified vanishing relative error
and the exponentially smaller target-relative error.
\end{proof}

This proposition is not merely a warning about rates.  The next theorem gives
two exact scalar sequences with the same source asymptotic and opposite
target behavior.

\section{Exact scalar underdetermination at balance}

Freeze the balance-phase clock
\begin{equation}\label{eq:sigma-balance}
 \sigma_k=\lambda^{-2(k-\eta_2)},
 \qquad
 C_*\sqrt{\sigma_k}=\frac2{C_M}\lambda^{-k}.
\end{equation}
For any desired scalar packet $0<X_k<r_H^{-2k}$, define
\begin{equation}\label{eq:inverse-parity}
 \delta_k(X)=1-\{1-r_H^{2k}X_k\}^{1/(2k)},
 \qquad
 \lambda_{-,k}(X)=-(1-\delta_k(X)).
\end{equation}
Then $\lambda_{-,k}(X)\in(-1,0)$ and its exact even parity packet is $X_k$.

\begin{theorem}[Opposite exact scalar completions]
\label{thm:scalar-completions}
For all sufficiently large $k$, let
\begin{equation}\label{eq:two-packets}
 X_k^{\rm close}=S_k,
 \qquad
 X_k^{\rm far}=S_k+\frac{\cA_{k,2k}}k.
\end{equation}
The two sequences \eqref{eq:inverse-parity} satisfy the same square-root law
\begin{equation}\label{eq:common-delta}
 \delta_k(X^{\rm close})
 =C_*\sqrt{\sigma_k}\{1+o(1)\},
 \qquad
 \delta_k(X^{\rm far})
 =C_*\sqrt{\sigma_k}\{1+o(1)\}.
\end{equation}
In the scalar information class with $Y_k=0$, the close sequence has
$p_{\sigma,k,2k}=0$, while the far sequence has
\begin{equation}\label{eq:far-residual}
 p_{\sigma,k,2k}=\frac{\cA_{k,2k}}k,
 \qquad
 \frac{|p_{\sigma,k,2k}|}{2H_k}
 =\frac1{C_Mk}(\beta R)^{2k}\{1+o(1)\}
 \longrightarrow\infty.
\end{equation}
\end{theorem}

\begin{proof}
By \eqref{eq:S-scale},
$r_H^{2k}X_k^{\rm close}=O(k\lambda^{-k})\to0$; the far correction is
smaller by a factor $O(k^{-1})$.  Hence both desired packets lie in the domain
of \eqref{eq:inverse-parity} eventually.  For $x_k=r_H^{2k}X_k\to0$,
\begin{equation}\label{eq:root-expansion}
 1-(1-x_k)^{1/(2k)}=\frac{x_k}{2k}\{1+o(1)\}.
\end{equation}
Equations \eqref{eq:A-scale}--\eqref{eq:S-scale} give, for both choices,
\[
 \frac{r_H^{2k}X_k}{2k}
 =\frac2{C_M}\lambda^{-k}\{1+o(1)\}
 =C_*\sqrt{\sigma_k}\{1+o(1)\},
\]
which proves \eqref{eq:common-delta}.  Substituting the two exact packets into
\eqref{eq:critical} with $Y_k=0$ gives zero and
$\cA_{k,2k}/k$.  Divide the latter by $2H_k$ and use
\eqref{eq:A-scale}.
\end{proof}

\begin{remark}[Strict information-class interpretation]
Equation \eqref{eq:inverse-parity} constructs scalar eigenvalue sequences,
not noisy transfer operators.  It proves that the exact parity-packet form
and the leading physical square-root law do not determine target closure at
the balance phase.  It does not replace the actual $\lambda_-(\sigma)$ or
construct two physical orbit-free remainders.
\end{remark}

\section{Protocol, verdict, and next physical route}

The executable artifact evaluates the exact inverse parity map at
$k=8,16,24,32$, verifies both packet recoveries, checks the common leading
square-root ratio, and records the growing far weighted residual.  Archived
decimals for $C_*$ and $C_M$ are reproduction inputs only, not interval
certificates or theorem evidence.

The scalar-only critical route is therefore \texttt{STOP\_SCOPED}: off the
unique double-alias phase it fails under a target-negligible actual remainder,
and at the balance phase the source scalar information admits opposite exact
target behaviors.  Actual signed critical compensation remains
\texttt{NOT\_TESTABLE}/open because $Y_k$ is not estimated.

The next physical layer is the order-$2k-2$ lower-sideband decomposition.  It
must retain the same noise clock, freeze the actual cells before evaluation,
and extract the complete period-$2(k-1)$ boundary orbit before testing any
scalar mechanism.  Closing one selected order would still not close the
remaining off-alias prefix \citep{WangSynchronizedPrefix2026}.

RH-288 remains inactive and the RH-341 frontier is not promoted
\citep{WangActualFrontier2026}.  Gates A--E remain false/open.  This paper
constructs no Hilbert--Polya operator, identifies no Riemann zero, proves no
von Mangoldt prime-power trace, proves no completed-zeta divisor equality,
and does not prove the Riemann Hypothesis.

\bibliographystyle{plainnat}
\bibliography{references}

\end{document}
